\chapter{Unit 4: Ideation and Creative Problem Solving Toolkits}
\label{chap:unit4-ideation-selection}

\section{The SCAMPER Creative Heuristic Matrix}

\begin{canvasbox}{Alex Osborn \& Bob Eberle's SCAMPER Framework}
\begin{tabularx}{\linewidth}{|l|X|X|} \hline
\textbf{Letter} & \textbf{Core Operative Trigger} & \textbf{Analytical Engineering Question} \\ \hline
\textbf{S - Substitute} & Replace materials, components, personnel, energy sources. & What material or software library can replace the current bottleneck? \\ \hline
\textbf{C - Combine} & Merge disparate functions, services, workflows, hardware. & Can we combine user authentication with payment checkout in 1 step? \\ \hline
\textbf{A - Adapt} & Import solutions from nature, biology, or other industries. & How does the airline industry solve standby overbooking? \\ \hline
\textbf{M - Modify / Magnify} & Exaggerate, scale up, shrink, change color, change cadence. & What happens if we magnify the interface 10x or reduce processing to 1s? \\ \hline
\textbf{P - Put to Other Use} & Repurpose residual outputs, waste energy, idle servers. & Can heat generated by servers warm office spaces? \\ \hline
\textbf{E - Eliminate} & Strip non-essential features, reduce complexity, remove forms. & What happens if we completely eliminate user passwords? \\ \hline
\textbf{R - Reverse / Rearrange} & Invert sequence, swap causes and effects, flip roles. & What if the customer grades the company before using the service? \\ \hline
\end{tabularx}
\end{canvasbox}

\section{Divergent Ideation Protocols}

\begin{frameworkbox}{1. Brainwriting 6-3-5 Protocol}
\begin{itemize}
    \item \textbf{Structure:} 6 participants, write 3 ideas each, passed around every 5 minutes for 6 rounds.
    \item \textbf{Yield:} Exactly $6 \times 3 \times 6 = 108$ distinct ideation variations generated in 30 minutes.
    \item \textbf{Pedagogical Advantage:} Eliminates loud-voice domination; guarantees equal cognitive contribution.
\end{itemize}
\end{frameworkbox}

\begin{frameworkbox}{2. Crazy Eights Divergence Sprint}
\textit{Fast sketching exercise where each participant folds a sheet of paper into 8 rectangles and sketches 8 distinct ideas in 8 minutes (1 minute per cell). Prevents self-censorship and forces breakthrough thinking.}
\end{frameworkbox}

\section{Solution Evaluation: Desirability, Feasibility, Viability (DFV)}

\begin{metricbox}{IDEO's Three Lenses of Innovation}
A viable solution exists at the central nexus where all three conditions are satisfied:
\begin{enumerate}[label=\textbf{\arabic*.}, leftmargin=2em]
    \item \textbf{Desirability (Human Lens):} Do users actually want, value, and emotionally embrace this solution?
    \item \textbf{Feasibility (Technical Lens):} Can it be engineered using accessible technology, architecture, and regulatory parameters?
    \item \textbf{Viability (Economic Lens):} Is the unit economics sustainable? Does it generate positive ROI or operational surplus?
\end{enumerate}
\end{metricbox}

\section{The Pugh Decision Matrix (Controlled Convergence)}

\begin{metricbox}{Stuart Pugh's Concept Selection Formula}
Select a known baseline (Concept $B$). Evaluate candidates across $k$ criteria:
\[
\text{Net Score} = \sum (+) - \sum (-)
\]
Where:
\begin{itemize}
    \item $(+)$ indicates candidate is superior to baseline on this criterion.
    \item $(-)$ indicates candidate is inferior to baseline on this criterion.
    \item $(0)$ indicates candidate performs equivalently to baseline.
\end{itemize}
Candidates with high negative net scores are eliminated; positive candidates are synthesized.
\end{metricbox}

\section{Prioritization \& Selection Matrices}

\begin{canvasbox}{The 2$\times$2 Impact vs. Effort Prioritization Matrix}
\begin{tabularx}{\linewidth}{|l|X|X|} \hline
 & \textbf{Low Effort (Fast / Cheap)} & \textbf{High Effort (Complex / Costly)} \\ \hline
\textbf{High Impact} & \textbf{Quick Wins:} Execute immediately. High payoff, low barrier. Builds team momentum. & \textbf{Strategic Bets / Major Projects:} Plan carefully, allocate senior engineers, decompose into iterative sprints. \\ \hline
\textbf{Low Impact} & \textbf{Fill-ins / Low-Hanging Fruit:} Execute during low-load intervals or assign as onboarding tasks. & \textbf{Money Pits / Thankless Tasks:} Avoid, deprioritize, or ruthlessly eliminate from roadmap. \\ \hline
\end{tabularx}
\end{canvasbox}

\begin{frameworkbox}{The MoSCoW Prioritization Rule}
\begin{itemize}
    \item \textbf{Must Have (M):} Non-negotiable regulatory, safety, or core functional requirements without which product cannot launch.
    \item \textbf{Should Have (S):} High-value features essential for competitive success, but temporary manual workarounds exist.
    \item \textbf{Could Have (C):} Desirable enhancements implemented only if surplus sprint capacity remains.
    \item \textbf{Won't Have (W):} Explicitly deferred out-of-scope features for current development horizon.
\end{itemize}
\end{frameworkbox}
